\usepackage[utf8]{inputenc}
\usepackage[T1]{fontenc}
\usepackage[ngerman]{babel}
\usepackage{lmodern}
\usepackage{microtype}
\usepackage{geometry}
\usepackage{booktabs}
\usepackage{longtable}
\usepackage{tabularx}
\usepackage{array}
\usepackage{enumitem}
\usepackage{csquotes}
\usepackage{xcolor}
\usepackage{graphicx}
\usepackage{tikz}
\usepackage{rotating}
\usepackage{hyperref}
\usepackage[backend=biber,style=ieee]{biblatex}

\usetikzlibrary{arrows.meta,positioning,fit,calc,shapes.geometric,shapes.misc,backgrounds,matrix}

\geometry{left=28mm,right=28mm,top=28mm,bottom=30mm}
\addbibresource{bibliography/references.bib}

\hypersetup{
  colorlinks=true,
  linkcolor=black,
  citecolor=black,
  urlcolor=blue,
  pdftitle={MBSE-gestuetzte Entwicklung eines mobilen Materialtransportroboters},
  pdfauthor={Projektteam}
}

\newcommand{\projectversion}{0.3}
\newcommand{\docstatus}{Arbeitsstand}
\newcommand{\sysml}{SysML v2}
\newcommand{\ros}{ROS2}

\newcolumntype{L}[1]{>{\raggedright\arraybackslash}p{#1}}
\newcolumntype{Y}{>{\raggedright\arraybackslash}X}

\setlist[itemize]{topsep=3pt,itemsep=2pt,parsep=0pt}
\setlist[enumerate]{topsep=3pt,itemsep=2pt,parsep=0pt}
\renewcommand{\arraystretch}{1.18}

\newenvironment{decisionbox}[1]
  {\par\medskip\noindent\textbf{#1}\par\smallskip\noindent}
  {\par\medskip}

% Modellgebundene Tabellenzeilen aus tables/generated/ (erzeugt mit tools/generate_docs.py).
% Das primitive \@@input funktioniert auch innerhalb von longtable.
\makeatletter
\newcommand{\modeltable}[1]{\@@input tables/generated/#1.tex }
\newcommand{\modeltext}[1]{\@@input tables/generated/#1.tex }
\makeatother
